% !TEX root = ../Presentation.tex
% ---------------------------------------------------------------------------
% Model comparison and recovery
% Reports: Model_B/Fit_B_Report.tex ("Comparison with Model A"),
%          Model_B/Model_Recov_BA_Report.tex, Model_C/Fit_C_Report.tex
%          ("Comparison with Model B")
% ---------------------------------------------------------------------------
\section{Model Comparison and Recovery}

\subsection{Predictive Performance}

\begin{frame}{Leave-one-out cross-validation}
  All three models fit to the same 12{,}937 trials in the same order, so the difference
  between two models can be computed trial by trial:
  \vspace{0.6em}

  \centering
  \begin{tabular}{lccc}
    \toprule
    Model & $\text{elpd}_\text{loo}$ (SE) & $p_\text{loo}$ & difference (SE) \\
    \midrule
    A: sessions pooled          & $-3344.0$ (64.5) & 13.4  & B $-$ A: $+164.1$ (19.2) \\
    B: sessions hierarchical    & $-3180.0$ (64.3) & 178.7 & \textbf{best} \\
    C: distractor as covariate  & $-3217.0$ (65.7) & 138.6 & C $-$ B: $-37.1$ (11.9) \\
    \bottomrule
  \end{tabular}

  \vspace{0.6em}
  \raggedright
  \begin{itemize}
    \item B beats A by 8.5 standard errors, despite 179 effective parameters against 13
    \item B beats C by 3.1 standard errors, although C needs fewer effective parameters
    \item Every Pareto $k$ below 0.7, so the PSIS-LOO estimates are reliable
  \end{itemize}
\end{frame}

\subsection{Model Recovery: B vs A}

% numbers from Model_B/Model_Recov_BA_Report.tex (made by Model_Recov_BA.py)
\begin{frame}{Model recovery: is B's advantage real?}
  Two questions before reading B's win as session-to-session variation:
  \begin{enumerate}
    \item Can LOO identify the model that generated a dataset of this size and design?
    \item Could B win just because it is more flexible? B contains A as the case of no
          between-session variation.
  \end{enumerate}
  \vspace{0.5em}
  Design:
  \begin{itemize}
    \item 10 replicates, each with two synthetic datasets on the real trial design
          (12{,}937 trials in 42 sessions)
    \begin{itemize}
      \item \textbf{data A:} responses from Model A, true values = A's real-data posterior means
      \item \textbf{data B:} session parameters drawn from B's real-data hyperparameters
    \end{itemize}
    \item Both models fit to every dataset; LOO picks the model with the higher
          $\text{elpd}_\text{loo}$; $\Delta = \text{elpd}(\text{B}) - \text{elpd}(\text{A})$
  \end{itemize}
\end{frame}

\begin{frame}{Model recovery: LOO picks the generating model}
  \centering
  \begin{tabular}{lcc}
    \toprule
     & selects A & selects B \\
    \midrule
    data from A & 9 & 1 \\
    data from B & 0 & 10 \\
    \bottomrule
  \end{tabular}

  \vspace{0.8em}
  \raggedright
  \begin{itemize}
    \item \textbf{Data from B:} B wins every time by a wide margin, $\Delta$ = 150.2 to 282.4
          ($z$ = 8.2 to 11.1); B's $p_\text{loo}$ = 167--190
    \item \textbf{Data from A:} A wins 9 of 10, but narrowly, $\Delta = -6.8$ to $+1.6$; the one
          B win is a near-tie ($z = 0.36$). B mimics A by shrinking its sds ($p_\text{loo}$ = 54--84)
    \item \textbf{Real data:} $\Delta = 164.1$ and $p_\text{loo} = 178.7$, inside the data-from-B
          range and far outside the data-from-A range
  \end{itemize}
\end{frame}

\plotframe{Model recovery: $\Delta$ in every replicate}
  {\plotsB/Model_Recov_Plots_BA/elpd_diff_by_replicate.png}

\begin{frame}{Model recovery: conclusions}
  \begin{itemize}
    \item With this design and data size, LOO identifies session-to-session variation
          reliably: B wins every time with $z > 8$ when the variation is there
    \item B's flexibility alone does not make it win: on data without session variation it
          loses or ties ($\Delta \le 1.6$)
    \item The real data look like the data-from-B replicates in both $\Delta$ and $p_\text{loo}$:
          the psychometric curves do vary between sessions
    \item The reverse, confirming that variation is absent, is only weakly answered: the more
          flexible model loses by just a few elpd points when the simpler one is true
    \item Limits: 10 replicates per model; one set of B hyperparameters (with smaller
          between-session sds the gap would shrink)
  \end{itemize}
\end{frame}

\subsection{Model C vs Model B}

% numbers from Model_C/Fit_C_Report.tex, "Comparison with Model B"
\begin{frame}{Model C vs Model B: where C loses}
  $\text{elpd}_\text{loo}(\text{C}) - \text{elpd}_\text{loo}(\text{B})$ summed over the trials at
  each amplitude; negative values favor Model B
  \vspace{0.2em}

  \centering\footnotesize
  \renewcommand{\arraystretch}{0.85}
  \begin{tabular}{rcccc}
    \toprule
     & \multicolumn{2}{c}{left} & \multicolumn{2}{c}{right} \\
    stimulus ($\mu m$) & bimanual & unimanual & bimanual & unimanual \\
    \midrule
    6  & $-6.6$ & $-3.5$ & $-5.1$ & $-2.3$ \\
    12 & $-4.2$ & $-3.0$ & $-1.8$ & $-2.8$ \\
    18 & 0.7    & $-1.6$ & $-3.4$ & $-0.3$ \\
    24 & $-1.0$ & 0.8    & 1.9    & 0.8 \\
    32 & $-0.1$ & $-3.4$ & 0.3    & 1.7 \\
    38 & $-0.3$ & $-1.4$ & 0.2    & 1.6 \\
    44 & 0.0    & $-3.8$ & 0.4    & $-1.5$ \\
    50 & $-0.1$ & 0.3    & 0.1    & 0.0 \\
    \bottomrule
  \end{tabular}

  \vspace{0.2em}
  \raggedright
  {\small
  \begin{itemize}
    \item 29.2 of the 37.1 are lost at 6 and 12~$\mu m$: one lapse rate for both conditions
    \item Next: left unimanual at 32 and 44~$\mu m$: one slope for both conditions
  \end{itemize}}
\end{frame}

\plotframe{Model C vs Model B: session PSE shifts}
  {\plotsC/Fit_Plots_C/pse_shift_C_vs_B.png}

\begin{frame}{Model C vs Model B: the weakest distractor}
  Session PSE shifts ($\mu m$): mean, and range of the posterior means, over the sessions at
  the weakest and strongest distractor
  \vspace{0.5em}

  \centering\small
  \begin{tabular}{llcc}
    \toprule
    hand & distractor & Model B & Model C \\
    \midrule
    left  & 6~$\mu m$ (8 sessions)   & $-1.13$ ($-3.82$ to 0.41)     & $-1.34$ ($-1.59$ to $-1.18$) \\
    left  & 32~$\mu m$ (13 sessions) & $-6.67$ ($-10.12$ to $-0.44$) & $-7.40$ ($-11.80$ to $-2.91$) \\
    right & 6~$\mu m$ (8 sessions)   & $-3.91$ ($-6.74$ to $-2.06$)  & $-1.75$ ($-1.88$ to $-1.64$) \\
    right & 32~$\mu m$ (13 sessions) & $-7.90$ ($-11.28$ to $-3.72$) & $-8.58$ ($-10.58$ to $-6.36$) \\
    \bottomrule
  \end{tabular}

  \vspace{0.5em}
  \raggedright
  \begin{itemize}
    \item The models agree on the average shift and, for the left hand, on the session shifts
          ($r = 0.91$)
    \item For the right hand at 6~$\mu m$, B's shift is more than twice C's: C ties the shift to
          the amplitude, so a weak distractor must give a small shift
    \item Suggests an effect of the distractor being present, on top of its amplitude
  \end{itemize}
\end{frame}

\begin{frame}{Model comparison: conclusions}
  \begin{itemize}
    \item \textbf{Sessions differ:} B predicts far better than A, and model recovery shows this
          is not B's flexibility alone.
    \item \textbf{The models agree on the PSE shift:} about $-4.2$ to $-4.5$~$\mu m$ (left) and
          $-5.7$ to $-6.1$~$\mu m$ (right) on average.
    \item \textbf{The shift depends on amplitude, but not only:} C confirms the trend inside a
          model, yet B still predicts better. C's misfit points to
    \begin{itemize}
      \item lapse rates per condition (the strongest evidence)
      \item a distractor-presence term, so the shift is not forced to 0 for a weak distractor
      \item possibly a bimanual slope, which would let the distractor change the JND
    \end{itemize}
    \item[$\Rightarrow$] Model B is the model to use for inference for now
  \end{itemize}
\end{frame}
